Inattention that is the same for rated and silent listings amounts to having fewer EPC-sensitive buyers. It scales both premia together and leaves the model of Section~\ref{sec:equilibrium} intact (part iv). Differential attention matters only if the default is optimistic (part iii); a buyer who, without thinking about energy, implicitly assumes an average dwelling is a cursed buyer by another name and stays below the ceiling.\footnote{If the default equals the prior mean cost $\bar\theta C_1+(1-\bar\theta)C_0$, it is at least $C_1$, and silence never outprices a verified efficient rating. If it equals the rational cost of the silent pool, attention is irrelevant. The evidence on shrouded costs cited above finds under-weighting toward zero rather than toward the mean, which supports an optimistic default. Housing-specific evidence on this point does not yet exist, and it is the assumption on which the extension most depends.} And the high-quality seller loses her automatic preference for disclosure once silence outprices a verified efficient rating, so that only the hiding cost stops her from hiding (part ii). Applied to the seven letters instead of the two quality levels, Lemma~\ref{lem:att_prices} also predicts where silence sits among ratings: the gap between a rated listing and a silent one should be linear in the letter's energy cost and cross zero at $C^*$. The model of Section~\ref{sec:equilibrium} places that crossing between D and E; limited attention with an optimistic default places it at or above A.
